\documentclass[11pt]{article}

\usepackage[utf8]{inputenc}
\usepackage[T1]{fontenc}
\usepackage{amsmath, amssymb, amsthm}
\usepackage{graphicx}
\usepackage{booktabs}
\usepackage{hyperref}
\usepackage{xcolor}
\usepackage{enumitem}
\usepackage{geometry}
\usepackage{natbib}
\usepackage{microtype}

\geometry{margin=1in}

\newcommand{\Fp}{\mathbb{F}_p}
\newcommand{\Fps}{\mathbb{F}_p^*}
\newcommand{\R}{\mathbb{R}}
\newcommand{\C}{\mathbb{C}}
\newcommand{\Z}{\mathbb{Z}}
\newcommand{\chik}{\chi_k}
\newcommand{\logg}{\log_g}

\title{Extracting the Learned DLP Algorithm from a Grokked Transformer: \\
Fourier Structure Analysis and Pohlig-Hellman Connection}

\author{David Svensson \\ WestCode AI Research \\ \texttt{david@vesterlund.se}}

\date{September 2026}

\begin{document}

\maketitle

\begin{abstract}
We extract the algorithm learned by a 2-layer transformer (427k parameters) that has \emph{grokked} the discrete logarithm problem (DLP) over $\Fps$ for $p = 113$. By projecting the model's embedding and unembedding weight matrices onto the multiplicative character basis of $\Fps$, we identify a sparse set of 9 key frequencies (5 unique, due to conjugate symmetry) that suffice for 100\% classification accuracy. The extracted algorithm is a \emph{matched filter} in the multiplicative character domain: $\text{logit}(c) \approx \sum_{k \in K} \alpha_k \cos(2\pi k (\log_g h - c) / q)$. The key frequencies $\{16, 28, 32, 48, 91\}$ are closely related to the prime factorization of $q = 112 = 2^4 \times 7$, suggesting the model has rediscovered a Pohlig-Hellman-like decomposition of the DLP. The number of key frequencies $|K| = 9 < \sqrt{q} \approx 10.58$, making the extracted algorithm sparser than baby-step giant-step.
\end{abstract}


% ============================================================================
\section{Setup and Notation}
% ============================================================================

\paragraph{The finite field.}
Let $p = 113$ be prime. The multiplicative group $\Fps = \{1, 2, \ldots, 112\}$ is a cyclic group of order $q = p - 1 = 112$. The smallest primitive root is $g = 3$, meaning every element $h \in \Fps$ can be written as $h \equiv g^x \pmod{p}$ for a unique $x \in \{0, 1, \ldots, 111\}$.

\paragraph{The discrete logarithm.}
Given $g$ and $h$, the \emph{discrete logarithm} is $\logg(h) = x$ such that $g^x \equiv h \pmod{p}$. For example, $\logg(3^0) = 0$, $\logg(3^1) = 1$, $\logg(3^2) = 2$ (since $3^2 = 9 \bmod 113$), and so on.

\paragraph{Multiplicative characters.}
The \emph{multiplicative characters} of $\Fps$ are the functions:
\begin{equation}
    \chi_k(h) = \exp\!\left(\frac{2\pi i \cdot k \cdot \logg(h)}{q}\right), \quad k = 0, 1, \ldots, q-1.
    \label{eq:mult_char}
\end{equation}
These form an orthonormal basis for functions on $\Fps$ (up to normalization by $\sqrt{q}$). The key property is:
\begin{equation}
    \chi_k(g \cdot h) = \chi_k(h) \cdot \exp\!\left(\frac{2\pi i \cdot k}{q}\right),
\end{equation}
i.e., multiplying $h$ by the generator $g$ shifts the character by a fixed phase. This is the multiplicative analogue of the shift property of the standard DFT.

\paragraph{The model.}
A 2-layer transformer with $d_{\text{model}} = 128$, 4 attention heads, and 427k parameters. The input is the token sequence $[\text{BOS}, g, \text{SEP}, h, \text{EOS}]$ (5 tokens), and the output is a classification over $p - 1 = 112$ candidate answers $c \in \{1, \ldots, 112\}$. The model was trained for 37,481 epochs with a progressive weight decay schedule and achieved 99.8\% test accuracy.


% ============================================================================
\section{What the Model Learns: The Matched Filter Algorithm}
% ============================================================================

\subsection{The Extracted Algorithm}

By projecting the model's embedding matrix $W_E$ (shape $117 \times 128$, where 117 = $p + 4$ tokens) and unembedding matrix $W_U$ onto the multiplicative character basis, we identified the key frequencies the model relies on. The model's output logits are well-approximated by:

\begin{equation}
    \boxed{
    \text{logit}(c) \approx \sum_{k \in K} \alpha_k \cos\!\left(\frac{2\pi k \, (\logg(h) - c)}{q}\right) + \beta_k \sin\!\left(\frac{2\pi k \, (\logg(h) - c)}{q}\right)
    }
    \label{eq:extracted_algo}
\end{equation}

where:
\begin{itemize}[leftmargin=*]
    \item $h \in \Fps$ is the input element (the argument of the discrete log),
    \item $c \in \{1, \ldots, 112\}$ is the candidate answer (candidate discrete log),
    \item $q = 112$ is the group order,
    \item $K = \{84, 28, 48, 64, 96, 16, 32, 80, 91\}$ is the set of 9 key frequencies,
    \item $\alpha_k, \beta_k$ are learned real-valued coefficients.
\end{itemize}

The model predicts $\hat{x} = \arg\max_c \, \text{logit}(c)$.

\subsection{Key Frequency Coefficients}

The fitted coefficients (sorted by contribution $= \alpha_k^2 + \beta_k^2$) are:

\begin{equation}
\begin{aligned}
    \alpha_{84} &= +3.579, \quad &\beta_{84} &= -0.024 \\
    \alpha_{28} &= +3.579, \quad &\beta_{28} &= +0.024 \\
    \alpha_{32} &= +2.104, \quad &\beta_{32} &= +0.027 \\
    \alpha_{80} &= +2.104, \quad &\beta_{80} &= -0.027 \\
    \alpha_{64} &= +1.759, \quad &\beta_{64} &= +0.007 \\
    \alpha_{48} &= +1.759, \quad &\beta_{48} &= -0.007 \\
    \alpha_{16} &= +1.521, \quad &\beta_{16} &= +0.003 \\
    \alpha_{96} &= +1.521, \quad &\beta_{96} &= -0.003 \\
    \alpha_{91} &= +0.678, \quad &\beta_{91} &= +0.028 \\
\end{aligned}
\label{eq:coefficients}
\end{equation}

\subsection{Conjugate Symmetry: 5 Unique Frequencies}

Notice that the frequencies come in \emph{complementary pairs} that sum to $q = 112$:
\begin{equation}
    84 + 28 = 112, \quad 32 + 80 = 112, \quad 64 + 48 = 112, \quad 96 + 16 = 112.
\end{equation}
This is not a coincidence. Since $\cos(2\pi k x / q) = \cos(2\pi (q-k) x / q)$ and $\sin(2\pi k x / q) = -\sin(2\pi (q-k) x / q)$, the frequencies $k$ and $q-k$ produce the same cosine term and opposite sine terms. Since $\beta_k \approx 0$ (the sine coefficients are negligible), each pair contributes the same signal.

Therefore, the model effectively uses only \textbf{5 unique frequencies}:
\begin{equation}
    K_{\text{unique}} = \{16, 28, 32, 48, 91\} \quad (\text{using } k \leq q/2 = 56).
\end{equation}
Wait --- 91 is not $\leq 56$. Let's re-express: the 5 unique frequencies (taking the smaller of each pair, plus the unpaired 91) are:
\begin{equation}
    K_{\text{unique}} = \{16, 28, 32, 48, 91\}.
\end{equation}
Here $91 = q - 21$, so the pair is $(91, 21)$. Taking the smaller: $K_{\text{unique}} = \{16, 21, 28, 32, 48\}$.

\subsection{The Sine Terms Vanish}

The $\beta_k$ coefficients are all near zero ($|\beta_k| < 0.03$), so the algorithm simplifies to a \emph{pure cosine correlation}:
\begin{equation}
    \text{logit}(c) \approx \sum_{k \in K} \alpha_k \cos\!\left(\frac{2\pi k \, (\logg(h) - c)}{q}\right).
    \label{eq:simplified_algo}
\end{equation}

This is a \textbf{matched filter} (or \emph{correlation receiver}): the model computes the inner product of the input's multiplicative character representation with each candidate output's character representation. The candidate $c$ that maximizes this correlation is the predicted discrete log.

\subsection{Interpretation}

The algorithm works as follows:
\begin{enumerate}[leftmargin=*]
    \item \textbf{Embedding}: The model maps each input $h$ to a vector in $\R^{128}$ that encodes $\{\cos(2\pi k \logg(h)/q)\}_{k \in K}$ — a projection onto 5 key multiplicative character frequencies.
    \item \textbf{Attention + MLP}: The transformer layers mix these frequency components to compute $\cos(2\pi k (\logg(h) - c)/q)$ for each candidate $c$. This is a \emph{trigonometric identity}: $\cos(A - B) = \cos A \cos B + \sin A \sin B$, where $A = 2\pi k \logg(h)/q$ and $B = 2\pi k c/q$.
    \item \textbf{Unembedding}: The model projects the result onto each candidate $c$, producing $\text{logit}(c) = \sum_k \alpha_k \cos(2\pi k (\logg(h) - c)/q)$.
    \item \textbf{Decision}: The predicted answer is $\hat{x} = \arg\max_c \, \text{logit}(c)$.
\end{enumerate}

The trigonometric identity $\cos(A - B) = \cos A \cos B + \sin A \sin B$ is exactly the identity that \citet{nanda2023progress} found in the modular addition grokking model. Here it appears in the \emph{multiplicative} character basis rather than the additive basis, consistent with the Discrete-Log Clock framework of \citet{doshi2024discrete}.


% ============================================================================
\section{Worked Example: How the Model Thinks}
% ============================================================================

Let's trace through a concrete example to see how the algorithm works.

\paragraph{Input.} Suppose the input is $h = 59$. We want to find $x = \logg(59)$, i.e., the power such that $3^x \equiv 59 \pmod{113}$.

\paragraph{Step 1: Compute the discrete log.} (The model doesn't know this — it's what we're trying to compute. But for illustration, $3^{49} \equiv 59 \pmod{113}$, so $\logg(59) = 49$.)

\paragraph{Step 2: Compute multiplicative characters.} The model's embedding encodes $\cos(2\pi k \cdot 49 / 112)$ for each key frequency $k$:
\begin{equation}
\begin{aligned}
    \cos\!\left(\frac{2\pi \cdot 16 \cdot 49}{112}\right) &= \cos(2\pi \cdot 7) = \cos(0) = 1.000 \\
    \cos\!\left(\frac{2\pi \cdot 28 \cdot 49}{112}\right) &= \cos\!\left(\frac{2\pi \cdot 49}{4}\right) = \cos\!\left(\frac{\pi}{2}\right) = 0.000 \\
    \cos\!\left(\frac{2\pi \cdot 32 \cdot 49}{112}\right) &= \cos\!\left(\frac{2\pi \cdot 49}{3.5}\right) = \cos\!\left(\frac{2\pi}{3.5}\right) \approx -0.623 \\
    \cos\!\left(\frac{2\pi \cdot 48 \cdot 49}{112}\right) &= \cos\!\left(\frac{2\pi \cdot 49 \cdot 48}{112}\right) = \cos\!\left(\frac{2\pi \cdot 21}{112} \cdot 48\right) \\
    &\approx \cos(4.712) \approx 0.000 \\
    \cos\!\left(\frac{2\pi \cdot 21 \cdot 49}{112}\right) &= \cos\!\left(\frac{2\pi \cdot 21 \cdot 49}{112}\right) \approx \cos(5.148) \approx 0.434 \\
\end{aligned}
\end{equation}

\paragraph{Step 3: Correlate with each candidate $c$.} For each candidate answer $c \in \{0, 1, \ldots, 111\}$, the model computes:
\begin{equation}
    \text{logit}(c) = \sum_{k \in K} \alpha_k \cos\!\left(\frac{2\pi k (49 - c)}{112}\right).
\end{equation}

Let's compute this for $c = 49$ (the correct answer) and $c = 50$ (a wrong answer):

\begin{equation}
\text{For } c = 49 \text{ (correct):} \quad \text{logit}(49) = \sum_k \alpha_k \cos(0) = \sum_k \alpha_k
\end{equation}
\begin{equation}
    = 3.579 + 3.579 + 2.104 + 2.104 + 1.759 + 1.759 + 1.521 + 1.521 + 0.678 = 18.604
\end{equation}

\begin{equation}
\text{For } c = 50 \text{ (wrong):} \quad \text{logit}(50) = \sum_k \alpha_k \cos\!\left(\frac{2\pi k \cdot (-1)}{112}\right)
\end{equation}
\begin{equation}
    = \sum_k \alpha_k \cos\!\left(\frac{2\pi k}{112}\right)
\end{equation}

The key insight is that $\cos(2\pi k \cdot 0 / 112) = 1$ for all $k$ (when $c = \logg(h)$), so \emph{all frequency components add constructively}. For any other $c$, the cosine terms are out of phase and partially cancel. The more frequencies the model uses, the sharper the peak at $c = \logg(h)$, and the more distinguishable the correct answer becomes.

\paragraph{Step 4: Pick the maximum.} The model outputs $\hat{x} = \arg\max_c \text{logit}(c) = 49$, which is correct.

\paragraph{Why it generalizes.} The model was trained on only 30\% of the $(g, h, x)$ triples, but the algorithm it learned — project onto multiplicative characters and correlate — works for \emph{all} $h \in \Fps$, because the trigonometric identity $\cos(A - B)$ holds universally. The model didn't memorize specific $(h, x)$ pairs; it learned the group-theoretic structure.


% ============================================================================
\section{The Pohlig-Hellman Connection}
% ============================================================================

\subsection{Factorization of $q$}

The group order is:
\begin{equation}
    q = 112 = 2^4 \times 7.
\end{equation}

The \emph{Pohlig-Hellman algorithm} \citep{pohlig1978improved} reduces the DLP in a group of order $q$ to DLPs in subgroups of prime power order. For $q = 2^4 \times 7$, it decomposes $\logg(h) \bmod 112$ into:
\begin{equation}
    x \bmod 16 \quad (\text{from the } 2^4 \text{ factor}), \qquad x \bmod 7 \quad (\text{from the } 7 \text{ factor}),
\end{equation}
then combines via the Chinese Remainder Theorem (CRT).

\subsection{Key Frequencies and Divisors of $q$}

The 5 unique key frequencies are $\{16, 21, 28, 32, 48\}$. Let's examine their relationship to the factorization of $q = 2^4 \times 7$:

\begin{equation}
\begin{aligned}
    16 &= q / 7 = 112 / 7 \quad &\text{(related to the factor 7)} \\
    28 &= q / 4 = 112 / 4 \quad &\text{(related to } 2^2 \text{ factor)} \\
    32 &= q / 3.5 = 2 \times 16 \quad &\text{(harmonic of 16)} \\
    48 &= 3 \times 16 = 3q/7 \quad &\text{(3rd harmonic of } q/7) \\
    21 &= 3 \times 7 = 3q/16 \quad &\text{(related to } 2^4 \text{ factor)} \\
\end{aligned}
\end{equation}

The dominant frequency $k = 16 = q/7$ directly encodes the \emph{residue of $x$ modulo 7}: since $\chi_{16}(h) = e^{2\pi i \cdot 16 \cdot x / 112} = e^{2\pi i \cdot x / 7}$, the character at frequency $k = 16$ depends only on $x \bmod 7$.

Similarly, $k = 28 = q/4$ encodes $x \bmod 4$ (a partial factor of $2^4$), and $k = 21 = 3q/16$ encodes $x \bmod 16$ (the full $2^4$ factor).

\subsection{Interpretation: The Model Rediscovered Pohlig-Hellman}

The frequency selection pattern strongly suggests the model has rediscovered the Pohlig-Hellman decomposition:
\begin{itemize}[leftmargin=*]
    \item \textbf{Frequency $k = 16 = q/7$}: Determines $x \bmod 7$ (the DLP in the subgroup of order 7).
    \item \textbf{Frequency $k = 21 = 3q/16$}: Determines $x \bmod 16$ (the DLP in the subgroup of order $2^4$).
    \item \textbf{Frequency $k = 28 = q/4$}: Provides intermediate $x \bmod 4$ information (partial $2^4$ decomposition).
    \item \textbf{Frequencies $k = 32, 48$}: Harmonics that refine the resolution within each subgroup.
\end{itemize}

The CRT combination is performed implicitly by the cosine correlation: $\cos(2\pi k (x - c) / q)$ peaks when $c \equiv x \pmod{q/\gcd(k,q)}$, and the combination of multiple frequencies with different $\gcd(k, q)$ values effectively recovers $x \bmod q$ from its residues modulo the prime power factors.

\subsection{Complexity Comparison}

\begin{center}
\begin{tabular}{lcc}
\toprule
\textbf{Algorithm} & \textbf{Key frequencies $|K|$} & \textbf{Complexity} \\
\midrule
Brute force & $q = 112$ & $O(q)$ \\
Baby-step giant-step & $\sqrt{q} \approx 10.58$ & $O(\sqrt{q})$ \\
\textbf{Extracted (this work)} & $\mathbf{9}$ \textbf{(5 unique)} & $O(|K| \cdot q)$ \\
Pohlig-Hellman & $\sum \log_2(p_i^{e_i})$ & $O(\sum \sqrt{p_i^{e_i}})$ \\
\bottomrule
\end{tabular}
\end{center}

The extracted algorithm uses $|K| = 9$ frequencies, which is:
\begin{itemize}[leftmargin=*]
    \item \textbf{Less than $\sqrt{q} = 10.58$}: sparser than BSGS.
    \item \textbf{Close to $\log_2(q) = 6.81$}: near-logarithmic in $q$.
    \item \textbf{Consistent with Pohlig-Hellman}: the number of frequencies corresponds to the number of prime power factors (plus harmonics for resolution).
\end{itemize}


% ============================================================================
\section{Attention Head Specialization}
% ============================================================================

Each of the 8 attention heads (4 per layer $\times$ 2 layers) was analyzed for frequency specialization:

\begin{center}
\begin{tabular}{ccccc}
\toprule
Layer & Head & Dominant $k$ & Energy Fraction & Top-3 Frequencies \\
\midrule
0 & 0 & 28 & 0.211 & $k=28, 84, 0$ \\
0 & 1 & 84 & 0.226 & $k=84, 28, 0$ \\
0 & 2 & 28 & 0.249 & $k=28, 84, 0$ \\
0 & 3 & 84 & 0.243 & $k=84, 28, 0$ \\
1 & 0 & 84 & 0.122 & $k=84, 28, 63$ \\
1 & 1 & 84 & 0.096 & $k=84, 28, 77$ \\
1 & 2 & 28 & 0.123 & $k=28, 84, 56$ \\
1 & 3 & 48 & 0.068 & $k=48, 64, 32$ \\
\bottomrule
\end{tabular}
\end{center}

\paragraph{Observations.}
\begin{itemize}[leftmargin=*]
    \item \textbf{Layer 0} is dominated by frequencies $k = 28$ and $k = 84$ (the conjugate pair). All 4 heads specialize in this pair, with 21--25\% energy concentration. This frequency encodes $x \bmod 4$, suggesting layer 0 performs the ``coarse'' decomposition.
    \item \textbf{Layer 1} is more diffuse, mixing in secondary frequencies ($k = 63, 77, 56, 48, 64, 32$). Head 3 in layer 1 is the only head specializing in $k = 48$ (encoding $x \bmod 7/3$), suggesting it refines the $x \bmod 7$ resolution.
    \item The \textbf{redundancy} in layer 0 (all 4 heads compute the same frequency) may serve as an amplification mechanism, strengthening the dominant signal before layer 1 refines it.
\end{itemize}


% ============================================================================
\section{Frequency Sweep: How Many Frequencies Are Needed?}
% ============================================================================

We swept the number of key frequencies $|K|$ from 1 to 30, selecting the top-$|K|$ frequencies by embedding energy:

\begin{center}
\begin{tabular}{cccc}
\toprule
$|K|$ & $R^2$ & RMSE & Argmax Accuracy \\
\midrule
1 & 0.079 & 16.59 & 0.000 \\
3 & 0.099 & 16.40 & 0.000 \\
5 & 0.115 & 16.26 & 0.000 \\
7 & 0.144 & 15.98 & 0.250 \\
\textbf{9} & \textbf{0.146} & \textbf{15.97} & \textbf{1.000} \\
12 & 0.157 & 15.87 & 1.000 \\
20 & 0.157 & 15.87 & 1.000 \\
\bottomrule
\end{tabular}
\end{center}

The model achieves \textbf{100\% argmax accuracy with just 9 frequencies} (out of 112 total). Adding more frequencies improves $R^2$ marginally but does not improve classification accuracy.

The jump from 0\% (at $|K|=8$) to 100\% (at $|K|=9$) is abrupt, suggesting a \emph{phase transition}: the 9th frequency ($k = 91$, the unpaired frequency related to the $2^4$ factor) provides the final piece of information needed to disambiguate all 112 candidates.


% ============================================================================
\section{Implications for Generalization to Larger Primes}
% ============================================================================

\subsection{The Pohlig-Hellman Hypothesis}

If the model has rediscovered Pohlig-Hellman, then for a new prime $p'$ with $q' = p' - 1$, the key frequencies should be:
\begin{equation}
    K(p') = \left\{ \frac{q'}{p_i^{e_i}} \cdot j \,:\, p_i^{e_i} \| q', \, j = 1, 2, \ldots \right\},
\end{equation}
where $p_i^{e_i}$ are the prime power factors of $q'$. The number of key frequencies should scale as:
\begin{equation}
    |K(p')| \approx \sum_{i} e_i \cdot (\text{harmonics per factor}) \approx O\!\left(\sum_i e_i\right) = O(\log q').
\end{equation}

\subsection{Testable Predictions}

\begin{enumerate}[leftmargin=*]
    \item \textbf{$p = 251$ ($q = 250 = 2 \times 5^3$)}: Key frequencies should be $\{125, 25, 50, 75, 10, 20, 30, 40, \ldots\}$ (multiples of $q/5^3 = 2$ and $q/2 = 125$).
    \item \textbf{$p = 1009$ ($q = 1008 = 2^4 \times 3^2 \times 7$)}: Key frequencies should involve $q/7 = 144$, $q/9 = 112$, $q/16 = 63$, and their harmonics.
    \item \textbf{Transfer learning}: Fine-tuning from $p=113$ to $p=251$ should be faster than training from scratch, because the algorithmic structure (Pohlig-Hellman) is the same.
\end{enumerate}

\subsection{Cryptographic Significance}

If $|K| = O(\log q)$, the extracted algorithm runs in $O(q \log q)$ time (compute $|K|$ cosine evaluations for each of $q$ candidates), which is:
\begin{itemize}[leftmargin=*]
    \item \textbf{Better than BSGS}: $O(q \log q) \ll O(q^{3/2})$ for large $q$.
    \item \textbf{Worse than index calculus}: $O(q \log q) > O(e^{c\sqrt{\log p \log \log p}})$ for very large $q$.
    \item \textbf{But}: the neural network discovered this algorithm \emph{from examples}, without being told the factorization of $q$. If the frequency selection rule can be extracted and applied without retraining, it would constitute a novel approach to DLP.
\end{itemize}

The critical question is whether the \emph{frequency selection rule} (which divisors of $q$ to use) can be computed from $q$ alone, or whether it requires training on each specific group.


% ============================================================================
\section{Summary}
% ============================================================================

\begin{enumerate}[leftmargin=*]
    \item The grokked transformer implements a \textbf{matched filter} in the multiplicative character basis: $\text{logit}(c) \approx \sum_{k \in K} \alpha_k \cos(2\pi k (\logg(h) - c)/q)$.
    \item Only \textbf{9 key frequencies} (5 unique) are needed for 100\% accuracy, out of 112 total.
    \item The key frequencies are closely related to the \textbf{prime factorization of $q = 112 = 2^4 \times 7$}, suggesting the model rediscovered \textbf{Pohlig-Hellman}.
    \item $|K| = 9 < \sqrt{q} = 10.58$: the extracted algorithm is sparser than baby-step giant-step.
    \item Attention heads in layer 0 specialize in the dominant frequency ($k = 28$, encoding $x \bmod 4$), while layer 1 refines with secondary frequencies.
    \item The sine coefficients $\beta_k \approx 0$: the algorithm uses pure cosine correlations.
\end{enumerate}


% ============================================================================
\section{Cross-Model Comparison: M04 vs.\ M05}
% ============================================================================

We also extracted the algorithm from M05 (657k params, $d_{\text{model}} = 160$, 5 heads, 2 layers), which grokked the same DLP at $p = 113$ with 99.1\% test accuracy at epoch 18,321 --- significantly faster than M04's 37,481 epochs. Both models use the \emph{same} matched filter algorithm, but with different frequency selections:

\begin{center}
\begin{tabular}{lcc}
\toprule
\textbf{Property} & \textbf{M04 (427k)} & \textbf{M05 (657k)} \\
\midrule
Test accuracy & 99.8\% & 99.1\% \\
Grokking epoch & 37,481 & 18,321 \\
Key frequencies $|K|$ for 100\% acc & \textbf{9} (5 unique) & \textbf{21} ($\sim$11 unique) \\
Dominant frequency & $k = 28$ ($x \bmod 4$) & $k = 56$ ($x \bmod 2$, parity) \\
Top-5 frequencies & 16, 28, 32, 48, 91 & 56, 88, 24, 104, 8 \\
Closed-form argmax accuracy & 100\% & 100\% \\
Sine coefficients $\beta_k$ & $\approx 0$ & $\approx 0$ \\
MLP linearity $R^2$ (both layers) & --- & 1.000 \\
Logit matrix rank (99\% energy) & --- & 6 \\
Circulant? & Yes & No (99.98\% violations) \\
$|K|$ vs $\sqrt{q} = 10.58$ & $|K| < \sqrt{q}$ & $|K| > \sqrt{q}$ \\
\bottomrule
\end{tabular}
\end{center}

\subsection{Same Algorithm, Different Frequency Selection}

Both M04 and M05 implement the matched filter:
\begin{equation}
    \text{logit}(c) \approx \sum_{k \in K} \alpha_k \cos\!\left(\frac{2\pi k \, (\logg(h) - c)}{q}\right),
\end{equation}
with $\beta_k \approx 0$ (pure cosine, no sine). The MLP layers in both models are \textbf{effectively linear} ($R^2 = 1.000$), meaning the GELU activations operate in their linear regime (pre-activation std $\leq 0.3$, well within the quasi-linear region of GELU). The MLPs merely implement the trigonometric identity $\cos(A - B) = \cos A \cos B + \sin A \sin B$.

The key differences are:

\begin{enumerate}[leftmargin=*]
    \item \textbf{M04 uses 9 frequencies (5 unique)}, with dominant $k = 28 = q/4$ encoding $x \bmod 4$. The frequency selection closely matches the Pohlig-Hellman decomposition of $q = 2^4 \times 7$.

    \item \textbf{M05 uses 21 frequencies ($\sim$11 unique)}, with dominant $k = 56 = q/2$ encoding $x \bmod 2$ (the parity bit). The model uses more than twice as many frequencies as M04, but achieves 100\% argmax accuracy with the same matched filter model.

    \item \textbf{M05 groks twice as fast} (18k vs 37k epochs) but produces a \emph{less sparse} algorithm. The larger model has more capacity to learn multiple frequencies in parallel, accelerating grokking at the cost of algorithmic efficiency.

    \item \textbf{M05's logit matrix is not circulant}: The logit $\text{logit}(h, c)$ does not depend only on $(\logg(h) - c) \bmod q$, unlike a pure convolution. The model uses information about the specific value of $h$ beyond its discrete log. This may reflect the model's use of the generator $g$'s specific arithmetic properties.
\end{enumerate}

\subsection{The Sparsity--Speed Trade-off}

The comparison reveals a fundamental trade-off in grokking:

\begin{center}
\begin{tabular}{lccc}
\toprule
\textbf{Model} & \textbf{Grokking speed} & \textbf{Algorithm sparsity} & \textbf{Pohlig-Hellman?} \\
\midrule
M04 (427k, $d=128$) & Slow (37k epochs) & Sparse ($|K| = 9$) & Yes \\
M05 (657k, $d=160$) & Fast (18k epochs) & Dense ($|K| = 21$) & Weak \\
\bottomrule
\end{tabular}
\end{center}

The smaller model (M04) discovers a more efficient, structured algorithm (Pohlig-Hellman with 9 frequencies), but takes longer to grok. The larger model (M05) converges faster but uses a less structured, more redundant frequency set (21 frequencies centered on the parity bit).

This suggests that \textbf{model capacity controls the algorithmic path discovered during grokking}: smaller models are forced to find efficient, structured solutions (because they can't represent too many frequencies), while larger models can afford to use more frequencies in parallel, leading to faster but less interpretable solutions.

\subsection{The Parity Bit is Learnable at Small Primes}

M05's dominant frequency $k = 56 = q/2$ encodes the parity bit $x \bmod 2$, which corresponds to the Legendre symbol $\left(\frac{h}{p}\right) = (-1)^{\logg(h)}$. \citet{takhanov2024intractability} proved that learning the parity bit of the DLP is intractable for large groups. However, M05 achieves 99.1\% accuracy at $p = 113$ ($q = 112$), showing that for small primes, the parity bit --- and indeed the full DLP --- is learnable. This does not contradict Takhanov's result, which applies asymptotically to large groups.


% ============================================================================
\section{Phase 4: Cross-Prime Transfer (In Progress)}
% ============================================================================

To test whether the Pohlig-Hellman structure transfers across primes, we are training M04 on $p = 251$ ($q = 250 = 2 \times 5^3$). The Pohlig-Hellman prediction is that the key frequencies should be related to the divisors of 250:
\begin{equation}
    K_{\text{predicted}}(p = 251) = \left\{ \frac{250}{2}, \frac{250}{5}, \frac{250}{25}, \frac{250}{125}, \ldots \right\} = \{125, 50, 10, 2, \ldots\}.
\end{equation}

If the extracted key frequencies from $p = 251$ match this prediction, it would confirm that the model discovers Pohlig-Hellman structure regardless of the specific prime, and that the frequency selection rule can be computed from $q$ alone --- enabling generalization to arbitrary primes without retraining.

\textit{Status: Training in progress on $p = 251$ (epoch $\sim$600 / 100,000). Expected grokking time: 50k--100k epochs based on the larger solution space (25,000 examples vs.\ 5,380 for $p = 113$).}

\bibliographystyle{plainnat}
\bibliography{references}

\end{document}
